%%%%%
%
%% Dans ce fichier, les �ventuels accents sont cod�s
%% en ISO-8859-15 (aka latin9)
%% La ligne suivante donne trois exemples 
%% e accent aigu : �, e accent grave : �, e-dans-l'o : �
\section{Correction de pics simples avec prise en compte des valeurs manquantes}

\itemstyle{Op\'erateur~:} \operfmt{Correc\_pics\_mnq}
\index{Correc\_pics\_mnq}

%%%%%OBJET%%%%
\itemstyle{Objet~:}
D�tection de pics sup�rieurs � un seuil se produisant sur une seule valeur. La valeur
"fausse" est �ventuellement remplac�e par la moyenne des valeurs qui l'encadrent en fonction du param�tre $N_r$ en entr�e. La d�tection de pics se fait aussi avant ou apr�s des valeurs manquantes.

%%%%%ARGUMENTS%%%%
\bigskip
%% Dans ce fichier, les �ventuels accents sont cod�s
%% en ISO-8859-15 (aka latin9)
%% La ligne suivante donne trois exemples 
%% e accent aigu : �, e accent grave : �, e-dans-l'o : �
%%%%%ENTREES%%%%
\begin{tabarg}
\multicolumn{3}{l}{\titretab {Entr\'ees}}\\
\hline
$X$ & Vecteur, \texttt{float} & Param�tre d'entr�e\\
$S_0$ & Coeff., \texttt{float} & Seuil de d�tection d'un pic\\
$S_{mnq}$ & Coeff., \texttt{float} & Seuil de d�tection d'un saut avant ou apr�s une valeur manquante \\
$N_r$ & Coeff., \texttt{float} & Nombre maximum de points invalid�s � remplacer par une interpolation lin�aire\\
\end{tabarg}

\medskip
%%%%%SORTIES%%%%
\begin{tabarg}
\multicolumn{3}{l}{\titretab {Sorties}}\\
\hline
$Flag$ & Vecteur, \texttt{float} & Flag de diagnostic de la correction des pics \\
$X_c$ & Vecteur, \texttt{float} & Param�tre corrig� (m�me unit� que $X$)\\
\end{tabarg}


%%%%%FORMULE%%%%
\itemstyle{Formule}

%% Formule de l'operateur : Correc_pics_mnq
%% Dans ce fichier, les �ventuels accents sont cod�s
%% en ISO-8859-15 (aka latin9)
%% La ligne suivante donne trois exemples 
%% e accent aigu : �, e accent grave : �, e-dans-l'o : �

A chaque instant $j$ :
\begin{itemize}
\item un pic est d\'etect\'e si : \newline
$X[j] - X[j-1] > S_0$ et $X[j] - X[j+1] > S_0$  (pic vers le haut)

ou 

$X[j-1] - X[j] > S_0$ et $X[j+1] - X[j] > S_0$  (pic vers le bas)

\hspace*{1cm} alors $flag[j] = 1$ si pic vers le haut, ou $flag[j] = -1$ si 
pic vers le bas \newline

Dans ce cas si le nombre de pics d�tect�s $\leq N_r$ alors $X_c[j]$ est remplac� par une interpolation lin�aire entre les valeurs valides, sinon $X_c[j]$ = mnq \newline

\item un saut apr�s une valeur manquante est d\'etect\'e si :\newline
$X[j-1] = mnq$
et
$|X[j+1] - X[j]| > S_{mnq}$ 

\hspace*{1cm} alors $flag[j] = 3$  si saut vers le haut, ou $flag[j] = -3$ si 
pic vers le bas 
et $X_c[j] = mnq$ \newline

\item un saut avant une valeur manquante est d\'etect\'e si :\newline
$X[j+1] = mnq$
et
$|X[j] - X[j-1]| > S_{mnq}$ 

\hspace*{1cm} alors $flag[j] = 2$ si saut vers le haut, ou $flag[j] = -2$ si 
saut vers le bas$X_c[j] = X[j]$ \newline

\item une forte variation positive est d\'etect\'e si :\newline
$X[j+1] - X[j] > S_0$ et $X[j] - X[j-1] > S_0$  (forte variation vers le haut)

\hspace*{1cm} alors $flag[j] = 4$ et  $X_c[j] = X[j]$ \newline

\item une forte variation n�gative est d\'etect\'e si :\newline
$X[j-1] - X[j] > S_0$ et $X[j] - X[j+1] > S_0$  (forte variation vers le bas)

\hspace*{1cm} alors $flag[j] = -4$ et $X_c[j] = X[j]$ \newline

\item aucune variation "particuli�re" n'est d�tect�e :\newline
\hspace*{1cm} alors $flag[j] = 0$ et $X_c[j] = X[j]$ \newline

\end{itemize}





%%%%%COMMENTAIRES%%%%
\itemstyle{Commentaires :}

Cet op�rateur remplace l'op�rateur  \operfmt{correc\_saut\_mnq}. \\
Il am�liore :
\begin{itemize}
\item  le traitements de tous les cas de sauts avant ou apr�s une valeur manquante
\item l'interpolation lin�aire en cas de pics successifs
\item la prise en compte du seuil de dur�e max ($N_r$) pour l'interpolation lin�aire ( il ne remplace pas les valeurs d�j� manquantes dans $X$)
\end{itemize}
\bigskip 

\emph{Attention} : L'interpolation lin�aire ne convient pas aux variables
qui peuvent �tre naturellement discontinues, tels les caps et
autres angles.





%%%%%AUTEUR%%%%
\itemstyle{Auteur~:}
CNRM/GMEI/TRAMM

%%%%%REFERENCE%%%%
\itemstyle{R\'ef\'erences~:}

